% ============================================================================
% Seccion 3: Ingenieria de Prompts y Adaptacion
% TFP - Rubrica: 2 puntos
% ============================================================================
\section{Ingenier\'ia de prompts y adaptaci\'on}

% TODO: Documentar las estrategias de prompting usadas:
% - System prompt para extraccion de restricciones estructuradas.
% - Chain-of-Thought para justificar el ranking.
% - Few-shot examples con respuestas canonicas.
% - Funcion de salida estructurada (JSON Schema / Pydantic) para los pesos.

% TODO: Documentar ajustes finos si se aplican (LoRA, adapters) o justificar
% por que se opto por prompting puro (costo, tiempo, datos disponibles).

% TODO: Describir la base vectorial: ChromaDB / FAISS / Pinecone,
% esquema de chunking, modelo de embeddings, estrategia de retrieval (top-k,
% MMR, reranking).

\subsection{Estrategia de prompting}

[Detallar los templates de prompt empleados:
- \texttt{system\_prompt\_extractor.py}: convierte el mensaje libre del usuario en un JSON estructurado (region, presupuesto, peso\_salario, peso\_costo, ...).
- \texttt{justification\_prompt.py}: recibe el ranking Top-N y redacta la explicacion en lenguaje natural con chain-of-thought.
- Few-shot examples: 3-5 ejemplos canonicos para anclar el formato.]

\begin{lstlisting}[style=pythonstyle, caption={System prompt para extraccion estructurada de preferencias}, label={lst:prompt-extract}]
EXTRACTOR_SYSTEM_PROMPT = """
Eres un asistente que convierte mensajes libres de un estudiante
peruano en preferencias estructuradas para un motor de recomendacion.

Devuelve un JSON con este esquema:
{
  "region": str | null,
  "presupuesto_max": float | null,
  "peso_salario": float,    # 0.0 - 1.0
  "peso_costo": float,
  "peso_admision": float,
  "peso_ubicacion": float
}

La suma de los cuatro pesos debe ser 1.0.
"""
\end{lstlisting}

\subsection{Flujo RAG (Retrieval-Augmented Generation)}

[Describir la base vectorial:
- Tecnologia: ChromaDB (local) o Pinecone (gestionado) [definir].
- Documentos indexados: descripciones de carreras, planes de estudio, testimonios.
- Modelo de embeddings: \texttt{sentence-transformers/all-MiniLM-L6-v2} o similar.
- Estrategia de retrieval: top-k=5, filtrado por metadata (region, area).
- Re-ranking: opcional con cross-encoder.]

\subsection{Ajustes finos (fine-tuning / LoRA)}

[Documentar si se realizo fine-tuning:
- Sobre que modelo base: [definir].
- Dataset de entrenamiento: tamano, origen.
- Hiperparametros: learning rate, epochs, LoRA rank.
- Justificacion si NO se realizo: costo, datos insuficientes, baseline suficiente.]